% !TEX program = xelatex
\documentclass[12pt,a4paper]{article}
\usepackage[UTF8]{ctex}
\usepackage{amsmath,amsfonts,amssymb}
\usepackage{graphicx}
\usepackage{booktabs}
\usepackage{hyperref}
\usepackage{geometry}
\geometry{margin=2.5cm}
\usepackage{caption}
\usepackage{float}

\title{\textbf{EmotionLens: Cross-Domain Facial Expression Recognition\\
with Controlled-Variable Experiments}}
\author{Group X \\ XMUM 2604 Natural Language Processing / Deep Learning}
\date{June 2026}

\begin{document}

\maketitle

\begin{abstract}
我们构建了一个跨域面部表情识别系统 EmotionLens。系统后端基于 ResNet18，训练数据融合了 FERPlus（10 人投票软标签）、RAF-DB（真实场景人脸）和自采集的 1,123 张多人脸数据。我们通过 6 组严格的控制变量实验，每次只改变一个因素，分别量化了架构（VGG16→ResNet18: −2.3pp）、标签质量（硬标签→10人投票软标签: +35.2pp）和数据多样性（单一源→多源融合: +20.0pp）对真实人脸泛化的各自贡献。实验表明，标签质量是最大的单一杠杆，数据多样性是恢复负面情绪识别的关键，而架构本身对真实场景泛化几乎无贡献。最优模型在 FERPlus 测试集上达到 89.0\%，在自采集真实人脸数据上达到 64.5\%，相比课程基线（VGG16+FER2013, 11.6\%）提升了 52.9 个百分点。模型被部署为一个 FastAPI + WebSocket 的实时 Web 应用，包含 5 个可切换的分析镜头，核心哲学是：引擎是可复用的核心，应用场景由用户自定义。
\end{abstract}

\tableofcontents
\newpage

%=====================================================================
\section{Introduction}
%=====================================================================

面部表情识别（Facial Expression Recognition, FER）是计算机视觉领域的一个重要任务，其目标是从人脸图像中自动识别出情感类别。基于 Ekman 和 Friesen (1971) 提出的七种基本情绪（neutral, happiness, surprise, sadness, anger, disgust, fear），FER 在课堂参与度监测、人机交互、在线监考、心理健康筛查等场景中具有广泛的应用潜力。

尽管现代深度网络在公开 FER 数据集上的准确率已超过 85\%（Barsoum et al., 2016; Li \& Deng, 2019），一个关键的瓶颈仍然存在：在实验室数据上训练出来的模型，面对真实场景下未见过的人脸时，表现通常会大幅下降。这一现象被称为"域间差距"（domain gap）或"主体身份偏差"（subject identity bias）——模型可能学会了识别特定人脸的 identity 特征，而不是表情本身的通用特征。

在本课程中，老师建议以 FER2013 数据集配合 VGG16 作为起点方案。我们首先完整复现了这一基线：在 FER2013 标准测试集上达到 69.5\%，但在我们自己拍摄的 1,123 张真实人脸数据上仅为 11.6\%——基本等同于随机猜测。这一巨大的性能落差成为了我们整个工作的出发点。

本文的工作围绕三个核心问题展开：（1）到底是什么因素决定了模型在真实人脸上的泛化能力？（2）能否通过严格的控制变量实验，将每个因素的贡献精确量化？（3）一个在真实人脸上可用的 FER 模型，能解锁哪些应用场景？

为此，我们设计了 6 组实验，每次只改变一个变量（架构、标签质量、数据多样性），构建了一条完整的因果链。在此基础上，我们将最优模型部署为一个实时 Web 应用 EmotionLens，包含 5 个可替换的分析镜头，用于展示核心引擎的通用性。

本文结构如下：Section 2 回顾相关工作；Section 3 介绍数据集；Section 4 详述模型实现和实验设计；Section 5 呈现结果和分析；Section 6 总结全文。

%=====================================================================
\section{Literature Review}
%=====================================================================

\subsection{FER 数据集的演进}

FER2013（Goodfellow et al., 2013）是首个大规模 FER 数据集，包含 35,887 张 48×48 灰度图像，由单标注者对每张图像赋予 0-6 的 hard label。然而，单一标注者的主观性和表情本身的模糊性导致 FER2013 存在显著的标注噪声。研究表明，FER2013 的标注一致性远低于多人投票方案（Barsoum et al., 2016）。

FERPlus（Barsoum et al., 2016）使用与 FER2013 完全相同的图像集，但改变了标注方式：每张图像由 10 人独立标注，最终以概率分布的形式呈现（例如，一张图像可能被标注为 [0.0, 0.1, 0.7, 0.1, 0.0, 0.0, 0.1]，表示 10 位标注者的投票分布）。这种软标签形式显著降低了标注噪声，同时新增了 contempt（蔑视）类别，形成 8 类体系。在我们的工作中，由于 RAF-DB 和自采集数据中 contempt 样本极少，我们丢弃该类别，统一为 7 类。

RAF-DB（Li \& Deng, 2019）是一个面向真实场景的 FER 数据集，包含 15,344 张高清人脸图像。与 FERPlus 不同，RAF-DB 的人脸经过对齐处理，标注质量更高，且图像来源覆盖更多样的真实环境。RAF-DB 在我们的实验中扮演了"域间桥梁"的角色——它弥补了 FERPlus 在真实场景多样性上的不足。

\subsection{架构选择}

VGG16（Simonyan \& Zisserman, 2014）是最经典的 CNN 架构之一，以堆叠的 3×3 卷积层和 138M 参数量著称。它在 ImageNet 上取得了里程碑式的成绩，但其巨大的参数量和相对较慢的推理速度使其在实时应用场景中不够理想。

ResNet（He et al., 2016）通过引入残差连接（residual connection）解决了深层网络的梯度消失问题，使得更深的网络也能稳定训练。ResNet18 仅有 11.2M 参数，约为 VGG16 的 1/12，同时得益于 ImageNet 预训练权重的迁移能力，在多个下游任务中表现出色。

EfficientNet（Tan \& Le, 2019）通过神经架构搜索（NAS）找到了深度、宽度和分辨率之间的最优平衡，在参数效率上表现优异。然而，在我们的初步实验中，EfficientNet-B0 在同等训练条件下仅达到 41.0\% 的 FERPlus 测试准确率，远低于 ResNet18 的 79.9\%，因此我们没有将其作为主力架构。

我们的选择并非基于"ResNet18 绝对最优"的假设，而是通过实际对比 VGG16、ResNet18 和 EfficientNet-B0 在相同数据和训练配方下的表现后得出的经验判断：ResNet18 在参数效率、训练稳定性和最终性能之间达到了最实用的平衡。

\subsection{跨域泛化与正则化}

域偏移（domain shift）是 FER 模型从实验室走向真实场景面临的核心挑战。训练分布（如 FERPlus 的 48×48 网上灰度图）与部署分布（如笔记本摄像头拍摄的室内人脸）之间的差异，会导致模型在未见域上的性能急剧下降（Chen et al., 2022）。

主体身份偏差（subject identity bias）是另一个关键问题：当模型在训练时见过某个人的部分表情，它可能学会通过识别"这个人的长相"而非"这个表情的特征"来进行分类。一旦面对全新的人脸，模型就会失效。这一现象在跨语料库 FER 研究中已被反复验证。

为缓解上述问题，我们采用了多种正则化技术：mixup（Zhang et al., 2018）通过在成对图像及其标签的凸组合上训练模型，增强了对域偏移的鲁棒性；label smoothing（Szegedy et al., 2016）将硬标签软化，防止模型过拟合虚假的确定性；AdamW（Loshchilov \& Hutter, 2019）通过解耦权重衰减机制，在保持 Adam 收敛速度的同时提供了更好的泛化性能。

%=====================================================================
\section{Data Set}
%=====================================================================

\subsection{公开数据集}

我们的训练语料来自三个公开数据源的融合：

\textbf{FERPlus}（Barsoum et al., 2016）：来自 Microsoft 官方发布，parquet 格式。共 35,803 张图像（train 28,709 / validation 3,546 / test 3,546），原始为 8 类软标签。我们丢弃 contempt 类别后统一为 7 类 Ekman 基本情绪。标签格式为 10 人投票的概率分布，在训练时使用 soft cross-entropy loss。

\textbf{RAF-DB}（Li \& Deng, 2019）：包含约 15,344 张真实场景人脸图像，分为 basic set（7 类）和 compound set（混合表情）。我们仅使用 basic set，标签为 hard label。RAF-DB 的人脸来自网络上的多样化场景，覆盖了比 FERPlus 更丰富的姿态、光照和人种。

\textbf{FER2013}（Goodfellow et al., 2013）：经典的 FER 基准数据集，CSV 格式，35,887 张 48×48 灰度图像。7 类 hard label。我们仅在 Exp1（VGG16 baseline）和 Exp1b（ResNet18 架构消融）中使用 FER2013，目的是与课程基线保持完全一致。

\subsection{自采集数据（Self-Captured）}

为评估模型在真实人脸——尤其是我们自己的脸——上的泛化能力，我们采集了一批自建数据。

\textbf{采集方式}：6 名组员各自使用笔记本电脑内置摄像头，在教室或宿舍室内自然光照条件下录制 7 种表情的视频片段。拍摄时，被摄者距离屏幕约 40-60cm，按提示依次做出 neutral、happiness、surprise、sadness、anger、disgust、fear 七种表情。

\textbf{标注流程}：录制完成后，脚本按固定间隔从视频中抽取帧。先用一个弱模型（小型预训练 FER 模型）对每帧进行预标注，提供粗略的表情预测作为参考。然后由组员逐帧人工审核，修正错误标签，删除模糊或无法判断的帧。最终保留 1,123 张经过审核的有效帧。

\textbf{数据分布}：6 名组员的贡献不均衡（图略，见 per\_user\_emotion\_bars.png），但总体上每种表情均有足够的样本覆盖。数据分割采用 stratified random split（80/20 train/val），而非按 subject 分割——这意味着同一人的不同帧可能同时出现在训练集和验证集中。我们在 Section 6 的 Limitations 中讨论了 subject-disjoint split 作为更好的替代方案。

\subsection{数据预处理}

所有图像在进入模型前经过统一的预处理管线：（1）调整为 224×224 像素 RGB 格式；（2）使用 ImageNet 均值和标准差进行归一化（$\mu=[0.485, 0.456, 0.406]$, $\sigma=[0.229, 0.224, 0.225]$）。对于 FER2013/FERPlus 的 48×48 灰度原图，先 resize 再复制到三个通道。RAF-DB 原始 RGB 图像直接 resize。自采集数据先经过 MTCNN（Zhang et al., 2016）人脸检测和裁剪，然后 resize 到统一尺寸。

\begin{table}[H]
\centering
\caption{数据统计总表}
\begin{tabular}{lcccc}
\toprule
数据源 & Train & Validation & Test & Total \\
\midrule
FER2013 & 28,709 & 3,589 & 3,589 & 35,887 \\
FERPlus & 28,709 & 3,546 & 3,546 & 35,803 \\
RAF-DB  & $\sim$12,000 & — & $\sim$3,000 & $\sim$15,344 \\
Self    & 899 & 224 & — & 1,123 \\
\bottomrule
\end{tabular}
\end{table}

%=====================================================================
\section{Implementation}
%=====================================================================

\subsection{Model Architecture}

我们选用 ResNet18（He et al., 2016）作为主要模型架构，使用 torchvision 提供的 ImageNet 预训练权重进行初始化。原始 ResNet18 的 1000-class 全连接层被替换为一个 Linear(512, 7) 层，将 512 维池化特征映射到 7 类情绪输出。输入尺寸为 224×224 RGB，使用 ImageNet 标准归一化。模型总参数量为 11,180,103。

选择 ResNet18 而非 VGG16 或 EfficientNet 的理由基于我们的实测：VGG16（138M 参数）在 FER2013 上仅达到 69.5\% 测试准确率，且推理速度较慢；EfficientNet-B0（5.3M 参数）在同等训练配方下仅达到 41.0\% FERPlus 测试准确率。ResNet18 在参数效率（11.2M）、训练稳定性和最终性能之间提供了最实用的折中。此外，残差连接机制有效缓解了深层网络的梯度消失问题，使得模型在较少的 epoch 内即可收敛。

\subsection{Training Recipe}

\subsubsection{Loss Function}

我们根据数据集的标签类型使用两种不同的损失函数：

\textbf{Hard Label Cross-Entropy}（用于 FER2013 和 RAF-DB）：
\begin{equation}
\mathcal{L}_{\text{hard}} = -\log p_{\text{model}}(y_{\text{true}})
\end{equation}
其中 $y_{\text{true}}$ 是单一的类别标签。这种 loss 强制模型以 100\% 置信度向该标签收敛，在标注存在噪声时可能导致模型学习到错误的确定性。

\textbf{Soft Cross-Entropy}（用于 FERPlus）：
\begin{equation}
\mathcal{L}_{\text{soft}} = -\sum_{k=1}^{K} p_{\text{voter}}(k) \cdot \log p_{\text{model}}(k)
\end{equation}
其中 $p_{\text{voter}}(k)$ 是 10 位标注者对类别 $k$ 的投票比例。这种 loss 允许模型学习一个与人类标注不确定性相匹配的概率分布，对标注噪声更加鲁棒。两种 loss 的对比本身也是实验设计的一部分——它直接量化了"标注质量"对泛化能力的影响。

\subsubsection{Optimizer and Gradient}

我们采用 AdamW 优化器（Loshchilov \& Hutter, 2019），其参数更新规则为：
\begin{equation}
\theta_{t+1} = \theta_t - \eta \cdot \left(\frac{\hat{m}_t}{\sqrt{\hat{v}_t} + \epsilon} + \lambda \theta_t\right)
\end{equation}
其中 $\hat{m}_t$ 和 $\hat{v}_t$ 分别为一阶和二阶动量的偏差修正估计，$\lambda$ 为解耦的权重衰减系数。相比标准 Adam，AdamW 将权重衰减从自适应学习率中解耦，在保持快速收敛的同时提供了更好的泛化性能。超参数设定为 $\beta_1=0.9$, $\beta_2=0.999$, $\eta=3\times10^{-4}$, $\lambda=1\times10^{-4}$。

训练使用混合精度（Automatic Mixed Precision, AMP），通过 torch.amp.GradScaler 实现：前向传播使用 FP16 以加速计算，反向传播时将 loss 乘以一个动态 scale factor 放大，防止小梯度在 FP16 下产生 underflow，unscale 后再执行 optimizer step。

\subsubsection{Learning Rate Schedule}

使用余弦退火调度（CosineAnnealingLR）：
\begin{equation}
\eta_t = \eta_{\min} + \frac{1}{2}(\eta_{\max} - \eta_{\min})\left(1 + \cos\left(\frac{t}{T_{\max}}\pi\right)\right)
\end{equation}
设定 $T_{\max}=35$ epochs，$\eta_{\max}=3\times10^{-4}$，$\eta_{\min}\approx 0$。选择余弦退火而非阶梯衰减的理由是：（1）无需手动设置衰减节点，减少了超参数调优的负担；（2）平滑的衰减曲线使得后期 epoch 可以在更精细的 loss 景观中探索，有助于找到更优的局部极小值。

\subsubsection{Sampling}

为解决 FER 数据集中固有的类别不平衡问题（如 happiness 样本远超 disgust），我们使用 WeightedRandomSampler。对每个类别 $c$，采样权重设定为：
\begin{equation}
w_c = \frac{1}{\text{count}(c)}
\end{equation}
采样时使用 replacement=True，每次采样数量等于数据集总样本数。这使得少数类（disgust, fear）在每个 epoch 中被多次采样，而多数类（happiness, neutral）的样本出现频率相应降低。这套采样器同时应用于 FER2013 和 FERPlus 的训练流程中，确保所有实验的采样策略完全一致。

\subsubsection{Regularization}

我们使用了三种互补的正则化技术：

\textbf{Label Smoothing}（$\varepsilon=0.1$）：将 hard label 的 one-hot 向量软化为 $y_{\text{smooth}} = (1-\varepsilon) \cdot y_{\text{hard}} + \varepsilon/K$，其中 $K=7$ 为类别数。这防止了模型对训练标签的过自信，提升了泛化能力（Szegedy et al., 2016）。

\textbf{mixup}（$\alpha=0.2$）：在每对训练样本 $(x_i, y_i)$ 和 $(x_j, y_j)$ 之间构建虚拟样本——$\tilde{x} = \lambda x_i + (1-\lambda) x_j$, $\tilde{y} = \lambda y_i + (1-\lambda) y_j$，其中 $\lambda \sim \text{Beta}(\alpha, \alpha)$。mixup 使模型在"插值"样本上训练，本质上是一种数据增强和正则化的结合，已被证明能有效提升模型对域偏移的鲁棒性（Zhang et al., 2018）。

\textbf{Weight Decay}（$1\times10^{-4}$）：通过 AdamW 的解耦机制施加，在每次参数更新时按比例缩小权重，阻止权重幅值无限制增长。

\subsubsection{Data Augmentation}

训练时对每张图像应用以下随机增强（概率独立）：RandomHorizontalFlip（$p=0.5$）、ColorJitter（brightness/contrast/saturation 各 0.2）、RandomErasing（$p=0.1$, scale=(0.02, 0.1)）。测试时使用水平翻转平均（H-flip averaging）作为 test-time augmentation：对每张图像分别用原图和水平翻转图进行推理，取平均概率作为最终预测。

\subsection{Experiment Design}

本章节是我们工作的核心方法论贡献。我们设计了 6 组实验，每次只改变一个变量，从而能够精确地将性能变化归因于特定的因素。

\begin{table}[H]
\centering
\caption{实验设计：控制变量矩阵}
\begin{tabular}{lccccl}
\toprule
实验 & 架构 & 训练数据 & 标签类型 & 控制变量 & 回答的问题 \\
\midrule
Exp1  & VGG16    & FER2013    & Hard & — (baseline)    & 老师方案表现如何？\\
Exp1b & ResNet18 & FER2013    & Hard & 仅架构          & 架构贡献了多少？\\
Exp2  & ResNet18 & FERPlus    & Soft & 仅标签质量      & 软标签提升了多少？\\
Exp3  & ResNet18 & F+R+S      & Mixed& 仅数据多样性    & 加数据提升了多少？\\
Exp4  & 全部4个  & —          & —    & 跨域评测        & 各模型的泛化能力？\\
Exp5  & FER库    & FER2013    & Hard & 第三方基线      & 现成方案表现如何？\\
\bottomrule
\end{tabular}
\end{table}

\textbf{统一的训练配方}：6 组实验使用完全相同的训练配置——AdamW optimizer（lr=$3\times10^{-4}$, wd=$1\times10^{-4}$）、CosineAnnealingLR（$T_{\max}=35$）、batch size=128、epochs=35（无 early stop, 跑满）。唯一例外是 Exp1 的 VGG16 因显存限制使用 batch size=64。这一统一性保证了当我们观察到性能差异时，差异的唯一起因就是我们有意改变的那个变量。

\textbf{控制变量逻辑链}：
\begin{enumerate}
    \item Exp1 $\to$ Exp1b：换架构（VGG16 $\to$ ResNet18），同数据（FER2013），同训练配方。
    \item Exp1b $\to$ Exp2：换标签质量（FER2013 hard labels $\to$ FERPlus soft labels），同架构（ResNet18）。
    \item Exp2 $\to$ Exp3：换数据多样性（FERPlus only $\to$ FERPlus + RAFDB + Self），同架构（ResNet18）。
\end{enumerate}

每一步只改变一个变量，整套实验构成了一个无懈可击的因果推理链条。

\subsection{Application Deployment}

我们将最优模型（ResNet18 + F+R+S）部署为一个实时 Web 应用 EmotionLens。后端使用 FastAPI + uvicorn + WebSocket：摄像头帧经 WebSocket 传至服务端，YuNet 人脸检测器定位人脸，ResNet18 推理得到 7 维情绪概率，经 EMA 时间平滑后返回前端。系统设计了 5 个独立的"分析镜头"，每个镜头消费同一条情绪流，将其加工为不同的指标和可视化：（1）群体情绪统计（示例场景：食堂满意度）；（2）突发情绪预警（示例场景：课堂异常检测）；（3）情绪时间线记录（示例场景：观影反应追踪）；（4）紧张度指标（示例场景：演讲教练）；（5）表情模仿游戏。前端使用原生 HTML/CSS/JS + Swiper.js 轮播。公网访问通过 Cloudflare Tunnel 提供的 HTTPS 隧道实现。

需强调的是，这 5 个镜头不是我们的产品——它们是核心引擎通用性的证明。引擎输出的是裸的 7 维概率分布，任何人都可以在不修改模型的情况下，编写自己的镜头来消费同一条情绪流。引擎是可复用的核心，镜头是开放式的示例。

%=====================================================================
\section{Results and Analysis}
%=====================================================================

\subsection{Control-Variable Chain}

图 \ref{fig:chain} 展示了控制变量链的完整结果。每一步只改变一个变量，Self data 上的准确率变化即为该变量的净贡献。

\begin{figure}[H]
\centering
\includegraphics[width=0.9\textwidth]{control_variable_chain.png}
\caption{控制变量链：从课程基线到最优模型的四步演进}
\label{fig:chain}
\end{figure}

\begin{itemize}
    \item VGG16 + FER2013: Self 11.6\% —— 课程基线，面对真实人脸基本失效。
    \item ResNet18 + FER2013: Self 9.2\% —— 换架构，不升反降（$-$2.3pp）。架构本身对真实泛化无贡献。
    \item ResNet18 + FERPlus: Self 44.4\% —— 换标签质量（hard $\to$ soft），单一最大提升（+35.2pp）。标注质量是第一个关键瓶颈。
    \item ResNet18 + F+R+S: Self 64.5\% —— 加数据多样性（+RAFDB+Self），再提升 +20.0pp。真实数据恢复了负面情绪的识别能力。
\end{itemize}

总提升幅度为 +52.9pp（11.6\% $\to$ 64.5\%）。

\subsection{Architecture Ablation}

表 \ref{tab:arch} 和图 \ref{fig:arch} 对比了 VGG16 和 ResNet18 在相同 FER2013 数据上的表现。

\begin{table}[H]
\centering
\caption{架构消融：VGG16 vs ResNet18（同数据、同训练配方）}
\label{tab:arch}
\begin{tabular}{lcc}
\toprule
指标 & VGG16 (138M) & ResNet18 (11M) \\
\midrule
FER2013 Test Acc & 69.46\% & 69.94\% \\
Self Data Acc    & 11.58\% & 9.24\% \\
\bottomrule
\end{tabular}
\end{table}

\begin{figure}[H]
\centering
\includegraphics[width=0.85\textwidth]{architecture_ablation.png}
\caption{架构消融：更先进的架构并未带来真实人脸性能的提升}
\label{fig:arch}
\end{figure}

在标准测试集上，两者的差距仅 0.48 个百分点；在真实人脸上，ResNet18 甚至比 VGG16 低了 2.34 个百分点。"换更好的架构就能提升真实场景性能"是一个直觉但不成立的假设——至少在本任务的范围内，架构选择的贡献微乎其微。

\subsection{Per-Class Recall Analysis}

图 \ref{fig:recall} 展示了全部 4 个模型在 Self data 上的 per-class recall。

\begin{figure}[H]
\centering
\includegraphics[width=\textwidth]{per_class_recall_all.png}
\caption{Per-class recall on self data：4 模型对比}
\label{fig:recall}
\end{figure}

FERPlus-only 模型在 sadness、disgust、fear 三个类别上的 recall 精确为 0\%——模型完全无法识别这些负面情绪。加入 RAFDB 和自采集数据后（F+R+S），这三个类别的 recall 分别恢复到 53.1\%、52.7\% 和 59.5\%。数据多样性的价值集中体现在负面情绪上：模型从未见过的 sad/disgust/fear 表情就是认不出来，只有加入了包含这些表情的真实场景数据后，模型才获得了辨别能力。

我们发现，经过多轮改进，我们的模型在 test 上的表现越来越好，但是在我们自建的数据集上表现还是与 FER 划分的测试集上的效果有显著差异。

而经过分析后，我们发现拆开 per-class recall 看，neutral (78\%)、surprise (70\%) 其实还不错，disgust (53\%)、anger (56\%)、fear (60\%) 表现却相对更加糟糕。我们推测这背后有两层原因。第一层，neutral 和 happiness 属于人类生物本能，表达方式跨文化高度一致——笑就是笑；但负面情绪的含义更抽象，不同的人理解和表达方式差异很大（Jack et al., 2012）。而另外一方面，FERPlus/RAF-DB 的标注以西方为主，但我们的 self test 全部来自亚洲组员，也许在西方人与亚洲人文化背景下对于这几种抽象表情的理解就有所差异，导致表现形式大相径庭。换句话说，这几类情绪的低 recall 可能不是模型没学好，而是标签一致性本身存在结构性的上限。

\begin{figure}[H]
\centering
\includegraphics[width=0.75\textwidth]{self_confusion_matrix.png}
\caption{最优模型（F+R+S）在 Self data 上的混淆矩阵}
\label{fig:cm}
\end{figure}

\subsection{Cross-Domain Evaluation}

表 \ref{tab:cross} 和图 \ref{fig:radar} 呈现了 4 个训练模型加 1 个 FER 库 baseline 在 3 个测试集上的完整评测结果。

\begin{table}[H]
\centering
\caption{跨域评测总表}
\label{tab:cross}
\begin{tabular}{lccc}
\toprule
模型 & FERPlus Test & FER2013 Test & Self Data \\
\midrule
VGG16 (FER2013)      & 7.1\%  & 69.5\% & 11.6\% \\
ResNet18 (FER2013)   & 9.0\%  & 69.9\% & 9.2\%  \\
ResNet18 (FERPlus)   & 79.9\% & 8.5\%  & 44.4\% \\
ResNet18 (F+R+S)     & \textbf{89.0\%} & 7.0\%  & \textbf{64.5\%} \\
EfficientNet-B0 (FERPlus) & 41.0\% & 10.6\% & 19.4\% \\
\midrule
FER Library          & —      & —      & 52.6\% \\
\bottomrule
\end{tabular}
\end{table}

\begin{figure}[H]
\centering
\includegraphics[width=0.8\textwidth]{cross_domain_radar.png}
\caption{跨域泛化雷达图：每个模型在自己训练的域上最优，仅多源模型跨域可用}
\label{fig:radar}
\end{figure}

每个模型在自己训练的领域上表现最好——这符合预期。但关键发现是：只有多源训练的 ResNet18 (F+R+S) 在真实人脸上达到了可用的性能水平（64.5\%）。FER 库（hustvl/fer，同样基于 FER2013 训练）在 self data 上为 52.6\%，但 disgust 召回率为 0\%，anger 仅为 1.8\%——这与 VGG16+FER2013 的表现模式一致，进一步印证了"FER2013 数据本身不足以训练出能泛化到真实人脸的模型"的结论。

\subsection{Domain Gap}

图 \ref{fig:gap} 展示了每个模型从实验室测试集到真实人脸之间的性能落差。

\begin{figure}[H]
\centering
\includegraphics[width=0.85\textwidth]{domain_gap.png}
\caption{Domain gap：Lab accuracy vs Self data accuracy}
\label{fig:gap}
\end{figure}

VGG16+FER2013 的域间差距高达 57.9pp，而 ResNet18+F+R+S 将差距缩小至 24.5pp。加入真实场景数据（RAFDB + Self）使域间差距缩小了一半以上。

\subsection{Contribution Decomposition}

图 \ref{fig:contrib} 将 52.9pp 的总提升分解为三个独立因素的贡献。

\begin{figure}[H]
\centering
\includegraphics[width=0.85\textwidth]{contribution_breakdown.png}
\caption{各因素对 Self data 准确率提升的贡献分解}
\label{fig:contrib}
\end{figure}

\begin{itemize}
    \item 架构（VGG16 $\to$ ResNet18）：$-$2.3pp —— 无贡献，甚至轻微负面。
    \item 标签质量（hard $\to$ 10-voter soft）：+35.2pp —— 最大的单一杠杆。
    \item 数据多样性（single $\to$ multi-source）：+20.0pp —— 第二关键因素，恢复负面情绪识别。
    \item 净提升：+52.9pp（11.6\% $\to$ 64.5\%）。
\end{itemize}

\subsection{FER Library Baseline}

我们使用 hustvl/fer 这一开源 FER 库对 self data 进行评估，作为第三方 baseline。该库同样基于 FER2013 训练。结果如下：总体准确率 52.6\%（看似不低），但 per-class recall 显示 disgust=0.0\%，anger=1.8\%，sadness=3.6\%——负面情绪几乎完全无法识别。这一表现模式与我们的 VGG16+FER2013（Exp1）高度一致，进一步验证了问题的根源在于 FER2013 数据本身的质量和多样性，而非我们的实现细节。

%=====================================================================
\section{Conclusion}
%=====================================================================

我们构建了一个跨域面部表情识别系统 EmotionLens，其核心是训练在 FERPlus 软标签、RAF-DB 真实人脸和自采集数据之上的 ResNet18 模型。通过 6 组严格的控制变量实验，我们将真实人脸准确率从课程基线的 11.6\% 提升至 64.5\%，并将这 52.9pp 的提升精确分解为三个独立因素的贡献：架构（$-$2.3pp，几乎无用）、标签质量（+35.2pp，最大的单一杠杆）和数据多样性（+20.0pp，恢复负面情绪识别的关键）。

本文的核心贡献不在于"达到了 64.5\%"这一数字本身，而在于我们通过控制变量链揭示了为什么——每一步只变一个变量，每一次变化都留下了清晰的因果足迹。这一方法论对任何构建 FER 系统的人都有参考价值：如果你只有一个改进预算，优先花在标签质量上，其次是数据的多样性，架构是最不值得投入的方向。

EmotionLens 的 5 个分析镜头（群体情绪统计、情绪预警、情绪时间线、演讲教练、模仿游戏）不是我们的产品，而是核心引擎通用性的证明。每个镜头消费同一条情绪流，无需修改模型。引擎是可复用的核心，镜头的设计空间由用户的想象力定义——课堂专注度监测、面试官反应分析、睡眠情绪日志……任何需要"从人脸读懂情绪"的场景，都可以在这个引擎上构建自己的镜头。

\section*{References}

\begin{enumerate}
    \item Barsoum, E., Zhang, C., Ferrer, C. C., \& Zhang, Z. (2016). Training deep networks for facial expression recognition with crowd-sourced label distribution. \textit{Proceedings of the 18th ACM International Conference on Multimodal Interaction (ICMI)}, 279--283.
    \item Chen, T., Pu, T., Wu, H., Xie, Y., Liu, L., \& Lin, L. (2022). Cross-domain facial expression recognition: A unified evaluation benchmark and adversarial graph learning. \textit{IEEE Transactions on Pattern Analysis and Machine Intelligence}, 44(12), 9887--9903.
    \item Ekman, P., \& Friesen, W. V. (1971). Constants across cultures in the face and emotion. \textit{Journal of Personality and Social Psychology}, 17(2), 124--129.
    \item Goodfellow, I. J., Erhan, D., Carrier, P. L., Courville, A., Mirza, M., Hamner, B., ... \& Bengio, Y. (2013). Challenges in representation learning: A report on three machine learning contests. \textit{International Conference on Neural Information Processing (ICONIP)}, 117--124.
    \item He, K., Zhang, X., Ren, S., \& Sun, J. (2016). Deep residual learning for image recognition. \textit{Proceedings of the IEEE Conference on Computer Vision and Pattern Recognition (CVPR)}, 770--778.
    \item Jack, R. E., Garrod, O. G. B., Yu, H., Caldara, R., \& Schyns, P. G. (2012). Facial expressions of emotion are not culturally universal. \textit{Proceedings of the National Academy of Sciences (PNAS)}, 109(19), 7241--7244.
    \item Li, S., \& Deng, W. (2019). Reliable crowdsourcing and deep locality-preserving learning for unconstrained facial expression recognition. \textit{IEEE Transactions on Image Processing}, 28(1), 356--370.
    \item Loshchilov, I., \& Hutter, F. (2019). Decoupled weight decay regularization. \textit{International Conference on Learning Representations (ICLR)}.
    \item Simonyan, K., \& Zisserman, A. (2014). Very deep convolutional networks for large-scale image recognition. \textit{arXiv preprint arXiv:1409.1556}.
    \item Szegedy, C., Vanhoucke, V., Ioffe, S., Shlens, J., \& Wojna, Z. (2016). Rethinking the inception architecture for computer vision. \textit{Proceedings of the IEEE Conference on Computer Vision and Pattern Recognition (CVPR)}, 2818--2826.
    \item Tan, M., \& Le, Q. V. (2019). EfficientNet: Rethinking model scaling for convolutional neural networks. \textit{Proceedings of the 36th International Conference on Machine Learning (ICML)}, 6105--6114.
    \item Zhang, H., Cisse, M., Dauphin, Y. N., \& Lopez-Paz, D. (2018). mixup: Beyond empirical risk minimization. \textit{International Conference on Learning Representations (ICLR)}.
    \item Zhang, K., Zhang, Z., Li, Z., \& Qiao, Y. (2016). Joint face detection and alignment using multitask cascaded convolutional networks. \textit{IEEE Signal Processing Letters}, 23(10), 1499--1503.
\end{enumerate}

\end{document}
